% ===== PRÉAMBULE CANONIQUE — à coller VERBATIM en tête de CHAQUE .tex =====
\documentclass[11pt,a4paper]{article}
\usepackage[utf8]{inputenc}\usepackage[T1]{fontenc}\usepackage{lmodern}
\usepackage[french]{babel}
\usepackage{amsmath,amssymb,mathtools}\usepackage{icomma}
\usepackage{siunitx}\sisetup{locale=FR}  % \qty{}{}, \si{}, \ang{}
\usepackage{xcolor}\usepackage{colortbl}
\usepackage{tikz}\usepackage{pgfplots}\pgfplotsset{compat=1.18}
\usepackage[siunitx,europeanresistors]{circuitikz} % circuits (élec.)
\usepackage[most]{tcolorbox}\usepackage{enumitem}\usepackage{array,booktabs}
\usepackage[a4paper,margin=2.2cm]{geometry}
\usetikzlibrary{arrows.meta,calc,angles,quotes,positioning,patterns,%
  backgrounds,shapes.geometric,decorations.pathmorphing,decorations.markings,intersections}
\definecolor{primary}{RGB}{222,49,99}\definecolor{primarydark}{RGB}{150,22,55}
\definecolor{defcol}{RGB}{0,114,178}\definecolor{methcol}{RGB}{52,92,148}
\definecolor{excol}{RGB}{0,135,90}\definecolor{attncol}{RGB}{200,40,55}
\tcbset{boxbase/.style={enhanced,breakable,boxrule=0.9pt,arc=2.5pt,
  left=3mm,right=3mm,top=2.3mm,bottom=2.3mm,fonttitle=\bfseries,coltitle=white}}
% --- boîtes TUTORIEL (noms EXACTS, ne pas renommer) ---
\newtcolorbox{cle}[1][]{boxbase,colback=defcol!6,colframe=defcol,
  title={\textsc{À retenir}\ifx\relax#1\relax\relax\else\textmd{ : \itshape #1}\fi}}
\newtcolorbox{methode}[1][]{boxbase,colback=methcol!7,colframe=methcol,sharp corners,
  title={\textsc{Méthode}\ifx\relax#1\relax\relax\else\textmd{ --- \itshape #1}\fi}}
\newtcolorbox{exemplebox}[1][]{boxbase,colback=excol!5,colframe=excol,
  title={\textsc{Exemple}\ifx\relax#1\relax\relax\else\textmd{ : \itshape #1}\fi}}
\newtcolorbox{piege}[1][]{boxbase,colback=attncol!6,colframe=attncol,
  title={\textsc{Piège}\ifx\relax#1\relax\relax\else\textmd{ : \itshape #1}\fi}}
\newtcolorbox{remarque}[1][]{boxbase,colback=black!4,colframe=black!45,
  title={\textsc{Remarque}\ifx\relax#1\relax\relax\else\textmd{ : \itshape #1}\fi}}
% --- boîtes EXERCICE / CORRIGÉ (noms EXACTS, auto-comptées) ---
\newtcolorbox[auto counter]{exo}[1][]{boxbase,colback=primary!4,colframe=primary,
  title={\textsc{Exercice}~\thetcbcounter\ifx\relax#1\relax\relax\else\textmd{ --- \itshape #1}\fi}}
\newtcolorbox[auto counter]{corrige}[1][]{boxbase,colback=excol!4,colframe=excol,
  title={\textsc{Corrigé}~\thetcbcounter\ifx\relax#1\relax\relax\else\textmd{ --- \itshape #1}\fi}}
\newcommand{\vv}[1]{\vec{#1}}\newcommand{\ds}{\displaystyle}
\newcommand{\R}{\mathbb{R}}\newcommand{\N}{\mathbb{N}}\newcommand{\Z}{\mathbb{Z}}
% =========================================================================
\begin{document}
\sisetup{per-mode=symbol}
% ================= DIFFICILE =================

\begin{exo}[masse ou poids ? --- vrai ou faux]
On emporte un même objet du niveau de la mer au sommet d'une montagne, puis sur la
Lune. Indique si chaque affirmation est \textbf{vraie} ou \textbf{fausse}.
\begin{enumerate}[label=\alph*)]
  \item Sa masse est la même au sommet de la montagne qu'au niveau de la mer.
  \item Son poids est légèrement plus faible au sommet (car $g$ y est un peu plus petit).
  \item Sur la Lune, sa masse est divisée par environ $6$.
  \item Sur la Lune, son poids est divisé par environ $6$.
\end{enumerate}
\end{exo}

\begin{exo}[deux masses en chute libre]
On lâche \emph{simultanément}, de la même hauteur et sans résistance de l'air, deux
objets de masses $m_1=\qty{2}{\kilogram}$ et $m_2=\qty{6}{\kilogram}$.
\begin{enumerate}[label=\alph*)]
  \item Écris la 2\up{e} loi de Newton pour un objet soumis à son seul poids, et
        montre que son accélération vaut $a=g$ (indépendante de la masse).
  \item Les deux objets touchent-ils le sol en même temps ? Justifie.
  \item Sur la Lune ($g$ plus faible), la chute serait-elle plus rapide ou plus lente ?
\end{enumerate}
\end{exo}

\begin{exo}[l'astronaute sur la Lune]
Un astronaute a une masse $m=\qty{90}{\kilogram}$. On prend
$g_{\text{Terre}}=\qty{9.8}{\metre\per\second\squared}$ et
$g_{\text{Lune}}=\qty{1.6}{\metre\per\second\squared}$.
\begin{enumerate}[label=\alph*)]
  \item Calcule son poids sur la Terre.
  \item Calcule son poids sur la Lune.
  \item Vérifie que le poids lunaire vaut environ le sixième du poids terrestre.
  \item Quelle est sa masse sur la Lune ?
\end{enumerate}
\end{exo}

\begin{exo}[abaisser un satellite : le piège de l'altitude]
Un satellite est à l'altitude $h=\qty{6400}{\kilo\metre}$. On l'abaisse à
$h'=\qty{3200}{\kilo\metre}$ (altitude divisée par $2$). On donne
$R_T=\qty{6400}{\kilo\metre}$.
\begin{enumerate}[label=\alph*)]
  \item Calcule les distances au centre $d=R_T+h$ et $d'=R_T+h'$. La distance
        au centre est-elle divisée par $2$ ?
  \item Par quel facteur la force gravitationnelle (donc $g$) est-elle multipliée ?
  \item La force est-elle multipliée par $4$, comme on pourrait le croire en
        « divisant l'altitude par 2 » ? Conclus.
\end{enumerate}
\end{exo}

\end{document}
